\subsection{Validación del prototipo}

La validación con personal médico y de enfermería no se ha realizado. El protocolo de usabilidad es una propuesta aún no ejecutada; no se dispone de participantes, observaciones, tiempos de tarea, errores observados ni puntajes SUS. Si se aplica, deberá distinguir la usabilidad percibida de la utilidad clínica~\cite{ref_brooke_sus}.

Como métodos de evaluación pendientes se consideran las siguientes actividades; su inclusión no afirma que se hayan ejecutado:

\begin{itemize}
    \item pruebas de los recorridos de consulta autorizada, notas, dictado, revisión de redacción y resumen de período;
    \item pruebas de integración con contratos de fuente comprometidos, usando datos sintéticos en la demostración;
    \item pruebas de autorización por permisos, institución y asignación asistencial aplicable;
    \item pruebas de ASR local en dispositivos representativos, registrando condiciones y mediciones;
    \item verificación del tratamiento local de audio y transcripción, y de la persistencia de notas confirmadas;
    \item pruebas de errores y recuperación que no impliquen operación offline completa;
    \item evaluación de usabilidad con personal asistencial, registrando método, participantes y resultados;
    \item revisión de resúmenes y propuestas con casos sintéticos, sin publicación ni escritura en la FCE.
\end{itemize}

Los resultados solo se informarán cuando exista evidencia reproducible de cada actividad. Las vistas desarrolladas pueden describirse como evidencia de implementación, pero no prueban usabilidad, eficacia, seguridad, calidad, latencia o aceptación del prototipo. Este apartado no afirma validación con participantes ni en producción, ni evaluación clínica.
